\documentclass[10pt,twocolumn]{article}
\usepackage[T1]{fontenc}
\usepackage[utf8]{inputenc}
\usepackage{lmodern}
\usepackage[margin=1.8cm,top=2.4cm,bottom=2cm,columnsep=0.6cm]{geometry}
\usepackage{fancyhdr}
\usepackage{titlesec}
\usepackage{booktabs}
\usepackage{amsmath,amssymb,amsfonts}
\usepackage{microtype}
\usepackage{xcolor}
\usepackage{mdframed}
\usepackage{parskip}
\usepackage{enumitem}
\usepackage{hyperref}

\definecolor{rulered}{RGB}{180,30,30}
\definecolor{boxgray}{RGB}{245,245,245}
\definecolor{keywordgray}{RGB}{100,100,100}

\pagestyle{fancy}
\fancyhf{}
\fancyhead[L]{\textsc{\small Claude Mag}}
\fancyhead[C]{\small\textit{Machine Learning Research Digest}}
\fancyhead[R]{\small 2026-07-13}
\fancyfoot[C]{\small\thepage}
\renewcommand{\headrulewidth}{0.8pt}

\titleformat{\section}{\normalsize\bfseries\color{rulered}}{}{0pt}{}%
  [\vspace{-4pt}\color{rulered}\rule{\columnwidth}{0.4pt}]
\titlespacing{\section}{0pt}{8pt}{4pt}

\hypersetup{colorlinks=true,urlcolor=rulered,linkcolor=black}

\begin{document}
\setlength{\parindent}{0pt}
\setlength{\parskip}{4pt}

{\small\color{keywordgray}\textbf{Keywords:}
  variational inference \textbullet{}
  continuous control \textbullet{}
  adaptive computation \textbullet{}
  reinforcement learning}\par\vspace{4pt}

{\LARGE\bfseries\raggedright
  Latent Memory Palace}\par\vspace{4pt}

{\large\itshape\raggedright
  Autoregressive Variational Inference Brings
  Chain-of-Thought Deliberation to Physical Control}\par\vspace{6pt}

{\small\textbf{By} Chuning Zhu, Eva Xu, Jose Barreiros,
  Krishnan Srinivasan, Paarth Shah \& Abhishek Gupta
  (Paul G.\ Allen School of CSE, University of Washington)}\par\vspace{2pt}

{\small\color{keywordgray}arXiv:2607.08724 \textbullet{}
  arXiv preprint, July 2026}
\par\vspace{2pt}\noindent{\color{rulered}\rule{\columnwidth}{0.6pt}}\vspace{6pt}

Language models reason flexibly by generating variable-length chains
of thought before producing an answer. Continuous control policies
cannot: they map observations directly to actions in a fixed-depth
computation graph, with no mechanism to think longer on harder problems.
Latent Memory Palace (LMP) closes this gap by formalizing deliberation
as autoregressive variational inference over a structured latent state
space, enabling physical agents to allocate more or fewer computation
steps depending on task complexity---achieving 16--26\% higher success
rates than fixed-depth recurrent baselines on long-horizon manipulation tasks.

\begin{mdframed}[backgroundcolor=boxgray,linecolor=rulered,linewidth=1pt,
    innerleftmargin=6pt,innerrightmargin=6pt,
    innertopmargin=6pt,innerbottommargin=6pt]
\textbf{\small AT A GLANCE}\par\vspace{4pt}
\begin{description}[leftmargin=0pt,labelsep=4pt,itemsep=2pt,parsep=0pt,
    font=\small\bfseries]
  \item[Problem:] Continuous control policies cannot adaptively allocate
    deliberation time; they use the same computation depth on easy and
    hard decisions.
  \item[Why it matters:] Long-horizon manipulation and dexterous tasks
    require proportionally deeper deliberation than routine motions,
    yet existing methods offer no mechanism for this.
  \item[Method:] Hierarchical variational autoencoder with autoregressive
    latent steps and a learned halting mechanism trained via ELBO
    maximization.
  \item[Result:] 47\% higher success rate on long-horizon manipulation
    versus best recurrent baseline (LSTM); 16--26\% across all tasks.
  \item[Evidence:] Five-task benchmark spanning stacking, tool use,
    pick-place, deformable manipulation, and multi-object coordination.
\end{description}
\end{mdframed}

\section{The Big Picture}

The gap between language model reasoning and physical control is
not merely architectural --- it is computational. A language model
generating a plan for a complex task can produce a chain of thought
of any length: simple tasks get short chains, complex tasks get long
ones. This adaptive computation is free in the discrete-token generation
paradigm where each token takes roughly the same wall-clock time.

Continuous control policies lack this flexibility entirely. A recurrent
policy (LSTM, GRU, S4) updates a fixed-size hidden state once per time
step. A diffusion policy iterates a fixed number of denoising steps.
Neither can elect to ``think harder'' about a particularly tricky
configuration. The result is policies that perform well on routine
subtasks but fail catastrophically when unexpected complexity arises.

Latent Memory Palace formalizes variable-depth deliberation in the
continuous domain by casting it as inference in a hierarchical latent
variable model where the number of latent computation steps is itself
a learned variable.

\section{Why Existing Approaches Fall Short}

\textbf{Recurrent policies} (LSTM, GRU, S4) maintain a hidden state
that accumulates information over time but applies the same fixed-depth
computation at each action step. The hidden state is informative but
the per-step computation is shallow.

\textbf{Model-based RL} (Dreamer, TDMPC) plans over fixed-depth horizon
rollouts. The horizon must be pre-specified and cannot adapt to task
difficulty without expensive online search.

\textbf{Diffusion policies} (Diffuser) perform iterative denoising in
action space. This iteration serves trajectory optimization, not
task-level deliberation --- the number of denoising steps is fixed and
does not correlate with decision complexity.

\textbf{Universal Transformers} adapt computation depth in sequence
modeling but have not been extended to continuous RL or equipped with
the variational training objective needed for learning under environment
uncertainty.

\section{Core Mechanism}

LMP introduces a \emph{latent deliberation sequence}: before producing
each action $a_t$, the agent runs $K_t$ autoregressive latent computation
steps $z_t^{(1)}, \ldots, z_t^{(K_t)}$. Each latent step refines the
agent's internal representation of the task state, analogous to one
token of chain-of-thought reasoning. The final latent state $z_t^{(K_t)}$
is then decoded to an action.

The number of steps $K_t$ is itself learned: a halt probability
$p_h(z_t^{(k)}, o_t) \in [0,1]$ is computed at each step, and the
agent halts when a Bernoulli sample fires. Harder observations
(estimated via an auxiliary complexity measure) produce lower halt
probabilities, naturally inducing more deliberation.

\section{Technical Formulation}

The model defines a hierarchical generative process with autoregressive
latent prior:
\[
p_\theta(z_t^{(1:K)} \mid o_t) = \prod_{k=1}^K
p_\theta\!\left(z_t^{(k)} \mid z_t^{(1:k-1)}, o_t\right)
\]

and posterior using retrospective information $h_t^{\text{ret}}$
(future observations and rewards, available during training):
\[
q_\phi(z_t^{(1:K)} \mid o_t, h_t^{\text{ret}}) = \prod_{k=1}^K
q_\phi\!\left(z_t^{(k)} \mid z_t^{(1:k-1)}, o_t, h_t^{\text{ret}}\right)
\]

Training maximizes the ELBO plus behavioral cloning on expert actions:
\[
\mathcal{L} = \mathbb{E}_{q_\phi}\!\left[
  \sum_t \log \pi_\psi(a_t \mid z_t^{(K_t)})
  - \beta \sum_t \mathrm{KL}\!\left(q_\phi \| p_\theta\right)_t
  - \lambda_c \sum_t K_t
\right]
\]

where the $\lambda_c \sum_t K_t$ term penalizes excess computation,
incentivizing economy in deliberation length.

\section{Learning or Inference Procedure}

\begin{enumerate}[itemsep=2pt,parsep=0pt]
  \item Collect offline demonstrations of target tasks.
  \item Train prior $p_\theta$, posterior $q_\phi$, policy $\pi_\psi$,
    and halt predictor $p_h$ jointly via ELBO on demonstration trajectories.
  \item At inference, use the prior $p_\theta$ only (no access to future):
    sample $z_t^{(k)} \sim p_\theta(\cdot \mid z_t^{(1:k-1)}, o_t)$
    and halt when $p_h(z_t^{(k)}, o_t) > u$ for $u \sim U[0,1]$.
  \item Decode the final latent state to action via $\pi_\psi(a_t \mid z_t^{(K_t)})$.
  \item Fine-tune with online RL (SAC) to improve beyond demonstrations.
\end{enumerate}

\section{What the Guarantee Says}

LMP does not provide formal sample-complexity bounds. The paper proves
empirically that: (1)~The average deliberation depth $\bar{K}$ correlates
positively with task difficulty (measured by expert-labeled complexity
scores, Pearson $\rho = 0.73$ across tasks). (2)~Performance improves
monotonically as the maximum allowed deliberation budget $K_{\max}$
increases, with diminishing returns above $K_{\max} = 8$ for the tested
tasks. (3)~The computation overhead of LMP versus a fixed-depth LSTM
is at most $2\times$ wall-clock time, as deliberation steps are shallow
and the learned halting mechanism allocates deep deliberation only to a
minority of high-complexity decisions.

\section{Experimental Findings}

\begin{center}
\small
\begin{tabular}{lcccc}
\toprule
\textbf{Task} & \textbf{LMP} & \textbf{LSTM} & \textbf{S4} & \textbf{Dreamer} \\
\midrule
Stacking (5 obj.)   & 71.3\% & 48.5\% & 52.1\% & 44.2\% \\
Multi-step tool use & 63.7\% & 41.2\% & 45.6\% & 38.9\% \\
Long-horizon P\&P   & 78.4\% & 56.8\% & 61.3\% & 52.7\% \\
Deformable manip.   & 55.1\% & 32.4\% & 37.8\% & 28.6\% \\
\bottomrule
\end{tabular}
\end{center}

Success rates (\%) averaged over 5 seeds and 200 evaluation episodes.
LMP exceeds the next-best baseline (S4) by 16--26\% across all tasks,
with the largest gains on deformable manipulation where unexpected
contact events require the most reactive deliberation.

\section{Ablations and Interpretation}

\textbf{Fixed vs.\ adaptive $K$:} Replacing the learned halting mechanism
with fixed $K = 4$ (the mean deliberation depth under adaptive halting)
reduces performance by 11\% on average, confirming that the \emph{alignment}
of depth with difficulty, not just depth on average, drives the gains.

\textbf{ELBO vs.\ BC-only:} Removing the ELBO training signal and
using behavioral cloning alone (no latent structure) reduces success
rates to near-LSTM levels, confirming the structured latent representation
is essential.

\textbf{Computation cost:} The deliberation overhead is 1.4--1.9$\times$
LSTM wall-clock time (task-dependent). The halting mechanism allocates
$K \geq 6$ steps to only 12\% of decisions, keeping average cost low.

\textbf{Transferability:} LMP policies pre-trained on one task set
transfer to related tasks with 40\% less fine-tuning data than LSTM
baselines, suggesting the latent memory structure captures more
transferable task representations.

\section*{Reference}

Chuning Zhu, Eva Xu, Jose Barreiros, Krishnan Srinivasan, Paarth Shah,
Abhishek Gupta.
\textbf{Latent Memory Palace: Reasoning for Control as Autoregressive
Variational Inference}.
arXiv:2607.08724, July 2026.
\url{https://arxiv.org/abs/2607.08724}

\end{document}
